\section{Vf-categories}

\subsection{The category of filters}

We begin by recalling the necessary background on \emph{filters}, so as to fix our notations and conventions. 

\todo{define the relative filter monad $\mathcal F \colon \Set \to \Pos$ and say that the notation $\phi_*(\tau)$ means $\mathcal F (\phi)(\tau)$. }

\begin{defn}
    The \emph{category of filters} $\FF$, introduced in \cite{koubekCategoryFilters1970}, is defined as follows:
    \begin{itemize}
        \item objects are pairs $(S, \sigma)$ of a set $S$ and a filter $\sigma$ on $S$;
        \item arrows $(T, \tau) \to (S, \sigma)$ are equivalence classes of functions $\phi \colon T \to S$ such that $\phi^{-1}(S_0) \in \tau $ for each $S_0 \in \sigma$, i.e.\ $\sigma \subseteq \phi_*(\tau)$, identified up to equality on a $\tau$-large subset.
    \end{itemize}
    We will often omit the carrier set of a filter, and we will not distinguish notationally between a function and the morphism it defines up to equivalence.
\end{defn}

We denote by $\delta_X$ the filter $\{X\}$ on a set $X$. The association $X \mapsto \delta_X$ realizes an embedding $\Set \hookrightarrow \FF$, which we will typically leave implicit. 

Recall from \cite[Thm.\ 6]{blassTwoClosedCategories1977} that an $\FF$-arrow $\phi \colon \tau \to \sigma$ is a monomorphism if and only if it is injective on a $\tau$-large subset, in which case it can be identified with the inclusion of a $\sigma$-large subset; on the other hand, it is an epimorphism if and only if $\phi(T_0) \in \sigma$ for each $T_0 \in \tau$, or equivalently $\sigma = \phi_*(\tau)$, in which case it can be assumed to be surjective. Moreover, as implied by the lemma below, $\FF$ is balanced.

\begin{lem}\label{lem:epis-are-regular}
    Every epimorphism in $\FF$ is effective. \todo{find ref in moerdijk and rephrase}
% \begin{proof}
%     Consider an epimorphism $\phi \colon \tau \epi \sigma$, i.e.\ such that $\sigma = \phi_*(\tau)$. Since $\FF$ is finitely complete (see, e.g., \cite[Thm.\ 9]{blassTwoClosedCategories1977}), it suffices to show that $\phi$ is the coequalizer of its kernel pair $\pi_1,\pi_2 \colon \tau\times_\sigma \tau  \rightrightarrows \tau$, so let $\chi \colon \tau \to \sigma'$ be another $\FF$-arrow such that $\chi \circ \pi_1 = \chi \circ \pi_2$. 
%     We define an $\FF$-arrow $\omega \colon \phi_*(\tau) \to \sigma'$ as follows. 
%     \[\begin{tikzcd}[row sep = 20pt]
% 	{\tau\times_\sigma\tau} & \tau & {\phi_*(\tau)} \\
% 	&& {\sigma'}
% 	\arrow["{\pi_1}", shift left, from=1-1, to=1-2]
% 	\arrow["{\pi_2}"', shift right, from=1-1, to=1-2]
% 	\arrow["\phi", two heads, from=1-2, to=1-3]
% 	\arrow["\chi"', from=1-2, to=2-3]
% 	\arrow["\omega", dashed, from=1-3, to=2-3]
%     \end{tikzcd}\]
%     For $\phi_*(\tau)$-almost-all $s\in S$ we know that $s = \phi(t)$ for some $t \in T$; thus, we set $\omega(s) = \chi(t)$. If $s = \phi(t')$ for some other $t' \in T$, then the pair $(t,t')$ lies in the carrier set of $\tau\times_\sigma \tau$, so that $\chi(t) = \chi \psi_1(t,t') = \chi \psi_2(t,t') = \chi(t')$. Moreover, $\sigma' \subseteq \omega_*(\phi_*(\tau))$: indeed, $\omega \circ \phi = \chi$ as functions, so that this reduces to $\sigma ' \subseteq \chi_* (\tau)$, which is true since $\chi$ is an $\FF$-arrow. Thus, $\omega$ is well-defined, and it is evidently the unique $\FF$-arrow such that $\omega \circ \phi = \chi$. 
% \end{proof}
\end{lem}

\begin{rmk}
    We thus also refer to epimorphisms in $\FF$ as \emph{covers}.
\end{rmk}

We denote by $0$ the initial object of $\FF$, isomorphic to any pair $(S, \mathcal P S)$. A filter $\sigma$ is \emph{proper} if the terminal map $\sigma \to 1$ is a cover, i.e.\ if every element of $\sigma$ is inhabited; clearly, this implies that (and is classically equivalent to) $\sigma$ is not isomorphic to $0$.

We recall here the \emph{dependent product} of filters, combining an indexed family of filters with respect to a filter on the index set. In the non-dependent case, this construction reduces to the tensor product of filters considered already since \cite{blassTwoClosedCategories1977}; in the case where the filter on the index set is actually a set, this recovers instead a description of coproducts in $\FF$. \todo{find a first ref for this too}

\begin{cons}[Dependent products]
    For a filter $\sigma$ on a set $S$ and a family $(\tau_s)_{s\in S}$ of filters, each on a set $T_s$, we denote by $\sigma \cdot \tau_s$ the filter on the set $\sum_{s\in S}T_s$ given by those subsets of the form $\sum_{s\in S_0} T_{s,0}$ where $S_0 \in \sigma$ and $T_{s,0} \in \tau_s$ for each $s \in S_0$. 
\begin{itemize}
    \item[(a)]In the case where all filters $\tau_s$ coincide with a single filter $\tau$ on a set $T$, the filter $\sigma\cdot \tau$ is the \emph{tensor product} of $\sigma$ and $\tau$, i.e.\ the filter on the set $S \times T$ consisting of those subsets $R$ such that 
    \[ \set{ s\in S | \set{ t \in T | (s,t) \in R } \in \tau } \in \sigma. \]
    \item[(b)]In the case where $\sigma$ is the filter $\delta_S$, the filter $S\cdot \tau_s$ is the \emph{coproduct} of the family $(\tau_s)_{s\in S}$, i.e.\ the filter on the set $\sum_{s\in S}T_s$ consisting of those subsets of the form $\sum_{s\in S} T_{s,0}$ where $T_{s,0} \in \tau_s$ for all $s \in S$. 
\end{itemize}
\end{cons}

Consider an $S$-family $(X_s)_{s\in S}$ of sets. Its \emph{reduced product} with respect to some filter $\sigma$ on $S$ is the set 
\[ \int_{s:\sigma} X_s \coloneqq \colim_{ S_0 \in \nu^{\op}} \prod_{s\in S_0} X_s,\]
whose elements are $\nu$-equivalence classes of tuples $(x_s \in X_s)_{s\in S_0}$ for some $S_0 \in \nu$, denoted as $(x_s)_{s:\sigma}$. In the case where all sets $X_s$ coincide with a single set $X$, we also speak of the \emph{reduced power} of $X$ with respect to $\sigma$, which we denote by $X^\sigma$. These constructions, well-known in model theory \cite{?}, generalize \emph{ultraproducts} and \emph{ultrapowers} to arbitrary filters. Using the category of filters $\FF$ we can describe them in a more useful way; let us start from reduced powers.

\begin{cons}[Reduced powers]\label{cons:reduced-powers}
For a set $X$ and a filter $\sigma$, the hom-set $\FF( \sigma, X)$ is in bijection with the reduced power $X^\sigma$. For an $\FF$-arrow $\phi \colon \tau \to \sigma$, the map $\FF(\sigma, X) \to \FF(\tau,X)$ given by precomposition with $\phi$ corresponds to the function defined by mapping $(x_{s})_{s:\sigma} \mapsto (x_{\phi t})_{t:\tau}$. These associations determine the \emph{reduced power functor}
\[ X^{-}\colon \FF^{\op} \to \Set.\]

% As explained in \cite{awodeyUltrasheavesDoubleNegation2004}, recall that we can consider the topology $J_\infty$ on $\FF$ generated by families $\set{ \gamma_i \colon \sigma_i \to \sigma }_{i\in I}$ such that $\bigcup_{i\in I} \gamma_i(S_i) \in \sigma$ for any $\set{ S_i \in \sigma_i }_{i\in I}$. As $\FF$ admits arbitrary (small) coproducts, it is immediate to see that this condition can be rephrased as asking for the family $\set{\gamma_i}$ to be jointly-epimorphic.

% \begin{lem}\label{lem:reduced-functors-are-sheaves}
%     For each set $X$, the functor $X^{-}\colon \FF^{\op} \to \Set$ is a $J_\infty$-sheaf.
%     \begin{proof}
%         It follows since the presheaf $X^{-}$ is represented by $\delta_X$ and the topology $J_\infty$ is subcanonical.
%     \end{proof}
% \end{lem}
\end{cons}

To generalize \Cref{cons:reduced-powers} to reduced products note first that, for an $S$-family $(X_s)_{s\in S}$ of sets, we can see $\sum_{s\in S} X_s$ as an object of $\FF/S$ via the projection. Moreover, filters on $S$ can be identified with objects of $\FF/S$ that are monomorphisms in $\FF$, which we can always assume to be represented by the identity map.

\begin{cons}[Reduced products]
   For an $S$-family $(X_s)_{s\in S}$ of sets and a filter $\sigma$ on $S$, the hom-set $\FF/S (\sigma, \sum_{s\in S} X_s)$ is in bijection with the reduced product $\int_{s:\sigma} X_s$.
\[\int_{s:\sigma} X_s \cong \left\{ \begin{tikzcd}[sep = small, cramped]
	& {\sum_{s\in S}X_s} \\
	\sigma & S
	\arrow[from=1-2, to=2-2]
	\arrow[dashed, from=2-1, to=1-2]
	\arrow[hook, from=2-1, to=2-2]
\end{tikzcd} \mbox{ in }\FF\right\} \]
As above, an $\FF$-arrow $\phi \colon \tau \to \sigma$ induces by precomposition a map $\FF/S(\sigma, \sum_{s\in S} X_s) \to \FF/S(\tau, \sum_{s\in S} X_s)$ corresponding to the function defined by mapping $(x_{s})_{s:\sigma} \mapsto (x_{\phi t})_{t:\tau}$. % and these associations determine the \emph{reduced product} functor
% \[ \int_{s:-} X_s \colon \FF^{\op} \to \Set.\]

\end{cons}

\UT{i'm not entirely sure that, in both constructions, we can say "corresponding to the function..." without assuming some degree of choice, in that we're defining a function on equivalence classes by looking at a representative}



\subsection{Vf-categories}

\begin{defn}
    A \emph{vf-category} $\X$ is defined by the following data:
    \begin{enumerate}
        \item a class $\X_0$ whose elements are called \emph{points};
        \item a class of \emph{vf-arrows}, whose elements have a \emph{domain} given by a point $p$, a \emph{codomain} given by a family of points $(q_s)_{s\in S}$ for some set $S$, and a \emph{type} given by a filter $\sigma$ on the set $S$, which we denote as $u \colon p \vuto{\sigma}q_s$; 
        \item for each point $p$, an \emph{identity} vf-arrow $\id_p \colon p \vuto{\ptuf} p$;
        \item for each vf-arrow $u \colon p\vuto{\sigma}q_s$ and each family $(v_s\colon q_s\vuto{\tau_s}r_{s,t})_{s\in S}$ of vf-arrows, a \emph{composition} $v_s\circ_\sigma u \colon p \vuto{\sigma\cdot\tau_s} r_{s,t}$, which we also denote graphically as
        \[v_s\circ_\sigma u \colon p \vuto{\sigma} (q_s\vuto{\tau_s} r_{s,t});\]
        \item\looseness=-1 for each vf-arrow $u \colon p \vuto{\sigma} q_s$ and each $\FF$-arrow $\phi \colon \tau \to \sigma$, a \emph{reindexing} $\phi^* u \colon p \vuto{\tau} q_{\phi t}$.
        
        % \item A vf-arrow (resp. vu-arrow) $\phi^*u : p\vuto{\tau}q_{\phi t}$ for each vf-arrow (resp. vu-arrow) $u : p\vuto{\sigma} q_s$ and each morphism of filters (resp. ultrafilters) $\phi : \tau\to\sigma$. This $\phi^*u$ is called the \emph{reindexing of $u$ along $\phi$}.
        % In the case $\phi$ injective (resp. surjective), we say that $\phi^*u$ is a \emph{specialization} (resp. a \emph{thickening}) of $u$.
        % In particular the reindexing of $\id_p$ along the unique map $\sigma\to\ptuf$ is called a \emph{diagonal} and is denoted by
        % \[\delta : a\vuto{\sigma} a.\]
        
    \end{enumerate}
    The above data must satisfy the following axioms:
    \begin{itemize}
        \item composition is unital and associative:
            \[u\circ_{\ptuf} \id_p = u\]
            \[\id_{q_s}\circ_{\sigma} u = u\]
            \[ (w_{s,t}\circ_{\tau_s} v_s)\circ_{\sigma} u = w_{s,t}\circ_{\sigma\cdot \tau_s} (v_s\circ_{\sigma} u)\]
        \item reindexing is functorial:
        \[\operatorname{id}^*_\sigma u = u\]
        \[\psi^*(\phi^*u) = (\phi\circ\psi)^*u\]
        \item reindexing is compatible with composition, in the sense that for all vf-arrows $u\colon p\vuto{\sigma} q_s$ and all families $(v_s \colon q_s\vuto{\tau_s} r_{s,t})_{s:\sigma}$ for some filter $\sigma$:
        \begin{enumerate}
            \item[] for each $\phi : \lambda\to\sigma$ with $\phi\cdot \operatorname{id}_{\tau_s} \colon \lambda\cdot\tau_{\phi \ell}\to\sigma\cdot\tau_s$
            \[v_{\phi \ell}\circ_{\lambda} (\phi^*u) = (\phi\cdot \operatorname{id}_{\tau_s})^*(v_s\circ_{\sigma} u),\]
            \item[] and for each $(\psi_s \colon \kappa_s\to\tau_s)_{s:\sigma}$ with $\operatorname{id}_{\sigma}\cdot\psi_s : \sigma\cdot\kappa_s\to\sigma\cdot\tau_s$
            \[({\psi_s}^*v_s)\circ_{\sigma} u = (\operatorname{id}_{\sigma}\cdot\psi_s)^*(v_s\circ_{\sigma} u).\] \UT{fix}
        \end{enumerate}
    \end{itemize}
    \todo{explain the notation in (vi) and (vii) and make the axioms slightly more readable}

We will typically write $\mathscr{X}$ in place of $\mathscr{X}_0$.
\end{defn}

For an $\FF$-arrow $\phi \colon\tau \to \sigma$, we also refer to the reindexing of $u \colon p\vuto{\sigma}q_s$ along $\phi$ as a \emph{specialization} of $u$ if $\phi$ is a monomorphism, or a \emph{thickening} of $u$ if $\phi$ is a cover. For each point $p$, we denote by $\diag_p^\sigma \colon p \vuto{\sigma} p$ the vf-arrow obtained by reindexing $\id_p \colon p\vuto{\ptuf} p$ along the unique $\FF$-arrow $!_\sigma \colon \sigma \to \ptuf$; we refer to vf-arrows of this form as \emph{diagonals}. 

\begin{defn}
    A vf-category $\X$ is \emph{locally small} if, for each point $p\in\X$, $S$-family of points $(q_s)_{s\in S}$, and filter $\sigma$ on $S$:
    \begin{enumerate}
        \item the class $\Hom_\sigma(p,q_s)$ of vf-arrows $p\vuto{\sigma}q_s$ is small, and 
        \item the presheaf 
        \[ \X_\sigma^{p,q_s} \colon (\FF/\sigma)^{\op} \to \Set \qquad (\phi \colon \tau\to\sigma) \mapsto \Hom_\tau(p,q_{\phi t})\]
        defined on arrows by reindexing is a \emph{small presheaf}, meaning that it is expressible as a small colimit of representables.
    \end{enumerate}
    \gabrielnote{Or equivalently, if this presheaf factorizes to $\Psh(\FF_\kappa/\sigma)\subseteq\Psh(\FF/\sigma)$ (resp. $\Psh(\UF_\kappa/\sigma)\subseteq\Psh(\UF/\sigma)$) for some large enough cardinal $\kappa$?} \UT{maybe we avoid speaking of cardinals if we want to stay constructive..} \gabrielnote{sure sure, here it would just mean a set}
    
    A vf-category $\X$ is \emph{small} if it is locally small and moreover its class of points is small, i.e.\ a set.
\end{defn}

\gabrielnote{rappeler les functors et nat transfo.}

We denote by $\vfCAT$ and $\vfCat$ the 2-categories of locally small and small vf-categories respectively. A \emph{vf-poset} is a vf-category such that for each point $p$, $S$-family of points $(q_s)_{s\in S}$, and filter $\sigma$ on $S$, the class $\Hom_\sigma(p,q_s)$ is a subsingleton; in case it is inhabited, we also write $p \vuleq{\sigma} q_s$. We denote by $\vfPOS$ and $\vfPos$ the full sub-2-categories of $\vfCAT$ and $\vfCat$ spanned by vf-posets.

\begin{ex}\label{ex:top-sp-as-vf-posets}
    A topological space $X$ can be canonically seen as a vf-poset by setting $p \vuleq{\sigma} q_s$ if every neighborhood $U \subseteq X$ of $p$ satisfies $q_s\in U$ for $s:\sigma$. \todo{in the preliminary, say that is a shorthand for  $(\forall s : \sigma) \ q_s \in U$, which in this case means that $q$ defines an element of $U^\sigma$} This definition gives rise to a 2-fully-faithful embedding $\Tinj \colon \TopSp \mono \vfPos$, where $\TopSp$ is regarded as a 2-category with the pointwise specialization order between continuous maps.
\end{ex}

\begin{ex}\label{ex:points-of-toposes}
    Let $\E$ be a topos. \todo{xx}
\end{ex}


We introduce here \emph{merging}, an operation on vf-arrows obtained combining diagonals and compositions. As shown by \Cref{lem:merging-is-inverse-to-specialization}, it provides an inverse operation for specialization along the canonical maps into a coproduct of filters.

\begin{cons}[Merging]
    For an $S$-family $(u_s \colon p \vuto{\tau_s} q_{s,t})_{s \in S}$ of vf-arrows with the same domain and a filter $\sigma$ on $S$, their \emph{merge} is the vf-arrow $p \vuto{\tau_i \cdot \sigma} q_{s,t}$ obtained as the composite
    \[ u_s \circ_\sigma \diag^\sigma_p \colon p \vuto{\sigma} (p \vuto{\tau_s} q_{s,t}),\]
    which we denote by $\bigmrg_{s:\sigma} u_s$. 
\end{cons}

\begin{lem}\label{lem:merging-is-inverse-to-specialization}
    Let $\mathscr{X}$ be a vf-category and fix a family $(\sigma_i)_{i\in I}$ of filters. For each point $p$ and each family $(q_{i,s})_{(i,s)\in \sum_i S_i}$ of points, the map 
    \[ \Hom_{\coprod_i \sigma_i} (p, q_{i,s} ) \to \prod_i \Hom_{\sigma_i} (p, q_{i,s})\]
    % \[\X(p, (q_{i,s})_{(i,s): \coprod_i \sigma_i } ) \to \prod_{i} \X(p, (q_{i,s})_{s:\sigma_i}) \]
    defined by specialization along each coprojection $\sigma_i \hookrightarrow \coprod_i \sigma_i$ is a bijection.
\begin{proof}
Fix a family of vf-arrows $(u_i \colon p \vuto{\sigma_i} q_{i,s})_{i\in I}$ and consider its merge $\bigmrg_{i:I} u_i$, which we can identify with a vf-arrow $p \vuto{\coprod_{i} \sigma_i} q_{i,s} $ since $I \cdot \sigma_i$ coincides with the coproduct $\coprod_i \sigma_i$. To see that this is indeed an inverse for the canonical map, consider any $i_0 \in I$; identifying the the coprojection $\iota_{i_0} \colon \sigma_{i_0} \mono \coprod_i \sigma_i$ with $i_0\cdot \operatorname{id} \colon 1 \cdot \sigma_{i_0} \to I \cdot \sigma_i$ where the $\FF$-arrow $i_0 \colon 1 \to I$ is represented by the function picking the element $i_0$ in $I$, \red{by the axioms} we have that 
    \[\textstyle \iota_{i_0}^* (\bigmrg_i u_i) = u_{i_0} \circ_1 (i_0^* \diag^I_p) = u_{i_0} \circ_1 \id_p = u_{i_0} .\]

On the other hand, consider a vf-arrow $u \colon p \vuto{\coprod_i \sigma_i} q_{i,s}$. Merging its specializations along each coprojection we obtain the vf-arrow $(\iota_i^* u)_i \circ_I \diag^I_p$. \red{By the axioms}, this coincides with the vf-arrow $(\operatorname{id}\cdot \iota_i )^* (u \circ_I \diag_p^I) $, where $\operatorname{id}\cdot \iota_i \colon I \cdot \sigma_i \to I\cdot (I\cdot \sigma_i)$. Moreover, again \red{by the axioms}, the vf-arrow $u\circ_I \diag^I_p = u \circ_I (!_I)^*\id_p$ coincides itself with $(!_I\cdot \operatorname{id} )^* u$, where $!_I \cdot \operatorname{id}_\sigma \colon I\cdot (I\cdot \sigma_i) \to \coprod_i \sigma_i$. Thus, since the composite 
\[\begin{tikzcd}
	{I\cdot \sigma_i} & {I \cdot (I\cdot \sigma_i)} & {I\cdot \sigma_i}
	\arrow["{\operatorname{id}\cdot \iota_i}", from=1-1, to=1-2]
	\arrow["{!_I\cdot\operatorname{id}}", from=1-2, to=1-3]
\end{tikzcd}\]
is the identity, we conclude:
\[ (\iota_i^* u)_i \circ_I \diag^I_p = (\operatorname{id}\cdot \iota_i )^* (u \circ_I \diag_p^I)  = (\operatorname{id}\cdot \iota_i )^* (!_I\cdot \operatorname{id} )^* u = u.\] 
\end{proof}
\end{lem}
